% GENERATED FILE. Edit canonical lessons, then run npm run generate:books.
% canonical-source-hash: fnv1a64:40432b4a539aeb96
% canonical-lessons: GU-C553-nagaru, GU-C553-sarnai, GU-C553-lalkarvu, GU-C553-ranakvu, GU-C553-thanganvu

\chapter{Drums and Pipes}
\label{ch:gu-drums-and-pipes}
\hlchaptermodality{553}

\begin{chapteropening}
\textbf{By the end of this chapter:} \emph{I can talk about loud instruments, bells and singing out.}

Complete the last lesson of chapter 553: I can talk about loud instruments, bells and singing out.
\end{chapteropening}

\section[\texorpdfstring{\gu{નગારું} (nagārũ)}{નગારું (nagārũ)}]{\gu{નગારું} (nagārũ) — a kettledrum}

\label{lesson:GU-C553-nagaru}



\begin{quote}
Before the new one: say the Gujarati for to make ring, then the Gujarati for to get going.
\end{quote}

\subsection*{You'll want to know: \gu{નગારું}}
\textbf{\gu{નગારું}} — \emph{nagārũ} — \textquotedblleft{}a kettledrum\textquotedblright{}.

\begin{quote}
\gu{જો} \gu{તહેવાર} \gu{હશે} \gu{તો} \gu{અરુણ} \gu{મંદિરમાં} \gu{નગારું} \gu{વગાડશે}.
\end{quote}

\begin{quote}
\emph{jo tahevār haśe to Aruṇ mandirmā̃ nagārũ vagāḍśe.}
\end{quote}

\begin{quote}
If there is a festival, Arun will play the kettledrum at the temple.
\end{quote}

A neuter noun.

\begin{cousinweb}
\textbf{\gu{નગારું}} (\emph{nagārũ}) is usually traced to Arabic \emph{naqqāra}, \textquotedblleft{}a kettledrum\textquotedblright{}, through Persian.
\end{cousinweb}

\begin{grammarlens}[title={What you've built}]
One more word for saying what happens if.
\end{grammarlens}

\subsection*{Your turn}
\begin{itemize}
\raggedright
  \item \emph{Say it:} \emph{nagārũ}
  \item \emph{Say it:} \emph{nagārũ}, once more
  \item \emph{Say it:} \emph{jo tahevār haśe to Aruṇ mandirmā̃ nagārũ vagāḍśe.} — If there is a festival, Arun will play the kettledrum at the temple.
  \item \emph{Recall:} say the Gujarati for to make ring, then the Gujarati for to get going, then say \emph{nagārũ} again
\end{itemize}

\begin{tcolorbox}[breakable,colback=teal!4,colframe=teal!35!black,title={Before you move on}]
What does \gu{નગારું} mean? (\textquotedblleft{}A kettledrum\textquotedblright{}.) Say it once more.
\end{tcolorbox}
\section[\texorpdfstring{\gu{શરણાઈ} (śarṇāī)}{શરણાઈ (śarṇāī)}]{\gu{શરણાઈ} (śarṇāī) — a shehnai (a reed pipe)}

\label{lesson:GU-C553-sarnai}



\begin{quote}
Before the new one: say the Gujarati for to get going, then the Gujarati for a kettledrum.
\end{quote}

\subsection*{You'll want to know: \gu{શરણાઈ}}
\textbf{\gu{શરણાઈ}} — \emph{śarṇāī} — \textquotedblleft{}a shehnai (a reed pipe)\textquotedblright{}.

\begin{quote}
\gu{જો} \gu{લગ્નમાં} \gu{શરણાઈ} \gu{નહીં} \gu{હોય} \gu{તો} \gu{મજા} \gu{નહીં} \gu{આવે}.
\end{quote}

\begin{quote}
\emph{jo lagnamā̃ śarṇāī nahī̃ hoy to majā nahī̃ āve.}
\end{quote}

\begin{quote}
If there's no shehnai at the wedding, it won't be any fun.
\end{quote}

A feminine noun.

\begin{cousinweb}
\textbf{\gu{શરણાઈ}} (\emph{śarṇāī}) is usually traced to Persian \emph{surnāy}, \textquotedblleft{}a festive reed pipe\textquotedblright{}.
\end{cousinweb}

\begin{grammarlens}[title={What you've built}]
One more word for saying what happens if.
\end{grammarlens}

\subsection*{Your turn}
\begin{itemize}
\raggedright
  \item \emph{Say it:} \emph{śarṇāī}
  \item \emph{Say it:} \emph{śarṇāī}, once more
  \item \emph{Say it:} \emph{jo lagnamā̃ śarṇāī nahī̃ hoy to majā nahī̃ āve.} — If there's no shehnai at the wedding, it won't be any fun.
  \item \emph{Recall:} say the Gujarati for to get going, then the Gujarati for a kettledrum, then say \emph{śarṇāī} again
\end{itemize}

\begin{tcolorbox}[breakable,colback=teal!4,colframe=teal!35!black,title={Before you move on}]
What does \gu{શરણાઈ} mean? (\textquotedblleft{}A shehnai (a reed pipe)\textquotedblright{}.) Say it once more.
\end{tcolorbox}
\section[\texorpdfstring{\gu{લલકારવું} (lalkārvũ)}{લલકારવું (lalkārvũ)}]{\gu{લલકારવું} (lalkārvũ) — to sing out}

\label{lesson:GU-C553-lalkarvu}



\begin{quote}
Before the new one: say the Gujarati for a kettledrum, then the Gujarati for a shehnai.
\end{quote}

\subsection*{You'll want to know: \gu{લલકારવું}}
\textbf{\gu{લલકારવું}} — \emph{lalkārvũ} — \textquotedblleft{}to sing out\textquotedblright{}.

\begin{quote}
\gu{જો} \gu{મીરા} \gu{ગીત} \gu{લલકારશે} \gu{તો} \gu{બધા} \gu{સાથે} \gu{ગાશે}.
\end{quote}

\begin{quote}
\emph{jo Mīrā gīt lalkārśe to badhā sāthe gāśe.}
\end{quote}

\begin{quote}
If Mira sings out a song, everyone will sing along.
\end{quote}

Past, with \gu{તેણે} (he or she …): \gu{તેણે} \gu{ગીત} \gu{લલકાર્યું}, \emph{teṇe gīt lalkāryũ}.

\begin{cousinweb}
\textbf{\gu{લલકારવું}} (\emph{lalkārvũ}) is built on \emph{lalkār}, \textquotedblleft{}a challenge, a shout\textquotedblright{}, which CDIAL files under \emph{*lallakka} (CDIAL 10973): to sing out loud and clear.
\end{cousinweb}

\begin{grammarlens}[title={What you've built}]
One more word for saying what happens if.
\end{grammarlens}

\subsection*{Your turn}
\begin{itemize}
\raggedright
  \item \emph{Say it:} \emph{lalkārvũ}
  \item \emph{Say it:} \emph{lalkārvũ}, once more
  \item \emph{Say it:} \emph{jo Mīrā gīt lalkārśe to badhā sāthe gāśe.} — If Mira sings out a song, everyone will sing along.
  \item \emph{Recall:} say the Gujarati for a kettledrum, then the Gujarati for a shehnai, then say \emph{lalkārvũ} again
\end{itemize}

\begin{tcolorbox}[breakable,colback=teal!4,colframe=teal!35!black,title={Before you move on}]
What does \gu{લલકારવું} mean? (\textquotedblleft{}To sing out\textquotedblright{}.) Say it once more.
\end{tcolorbox}
\section[\texorpdfstring{\gu{રણકવું} (raṇakvũ)}{રણકવું (raṇakvũ)}]{\gu{રણકવું} (raṇakvũ) — to chime}

\label{lesson:GU-C553-ranakvu}



\begin{quote}
Before the new one: say the Gujarati for a shehnai, then the Gujarati for to sing out.
\end{quote}

\subsection*{You'll want to know: \gu{રણકવું}}
\textbf{\gu{રણકવું}} — \emph{raṇakvũ} — \textquotedblleft{}to chime\textquotedblright{}.

\begin{quote}
\gu{જો} \gu{ઘંટડી} \gu{રણકશે} \gu{તો} \gu{તરત} \gu{બારણું} \gu{ખોલો}.
\end{quote}

\begin{quote}
\emph{jo ghanṭaḍī raṇakśe to tarat bārṇũ kholo.}
\end{quote}

\begin{quote}
If the bell rings, open the door at once.
\end{quote}

Past (it …): \gu{તે} \gu{ઘંટડી} \gu{રણકી}, \emph{te ghanṭaḍī raṇkī}.

\begin{cousinweb}
\textbf{\gu{રણકવું}} (\emph{raṇakvũ}) is imitative, the ring of a bell or of coins in the word itself. CDIAL lists Gujarati \emph{raṇko}, \textquotedblleft{}a tinkling\textquotedblright{}, and \emph{raṇraṇvũ}, \textquotedblleft{}to tinkle\textquotedblright{}, under Sanskrit \emph{raṇa}, \textquotedblleft{}noise\textquotedblright{} (CDIAL 10596); the verb itself is not listed there.
\end{cousinweb}

\begin{grammarlens}[title={What you've built}]
One more word for saying what happens if.
\end{grammarlens}

\subsection*{Your turn}
\begin{itemize}
\raggedright
  \item \emph{Say it:} \emph{raṇakvũ}
  \item \emph{Say it:} \emph{raṇakvũ}, once more
  \item \emph{Say it:} \emph{jo ghanṭaḍī raṇakśe to tarat bārṇũ kholo.} — If the bell rings, open the door at once.
  \item \emph{Recall:} say the Gujarati for a shehnai, then the Gujarati for to sing out, then say \emph{raṇakvũ} again
\end{itemize}

\begin{tcolorbox}[breakable,colback=teal!4,colframe=teal!35!black,title={Before you move on}]
What does \gu{રણકવું} mean? (\textquotedblleft{}To chime\textquotedblright{}.) Say it once more.
\end{tcolorbox}
\section[\texorpdfstring{\gu{થનગનવું} (thanganvũ)}{થનગનવું (thanganvũ)}]{\gu{થનગનવું} (thanganvũ) — to be itching to dance}

\label{lesson:GU-C553-thanganvu}



\begin{quote}
Before the new one: say the Gujarati for to sing out, then the Gujarati for to chime.
\end{quote}

\subsection*{You'll want to know: \gu{થનગનવું}}
\textbf{\gu{થનગનવું}} — \emph{thanganvũ} — \textquotedblleft{}to be itching to dance\textquotedblright{}.

\begin{quote}
\gu{જો} \gu{ગરબાનું} \gu{સંગીત} \gu{શરૂ} \gu{થશે} \gu{તો} \gu{બાળકો} \gu{થનગનશે}.
\end{quote}

\begin{quote}
\emph{jo garbānũ sangīt śarū thaśe to bāḷko thanganśe.}
\end{quote}

\begin{quote}
If the garba music starts, the children will be itching to dance.
\end{quote}

Past (he …): \gu{તે} \gu{થનગન્યો}, \emph{te thanganyo}.

\begin{cousinweb}
\textbf{\gu{થનગનવું}} (\emph{thanganvũ}) is imitative, the tapping of eager feet. No dictionary we use records an older form, so its history is not settled.
\end{cousinweb}

\begin{grammarlens}[title={What you've built}]
One more word for saying what happens if.
\end{grammarlens}

\subsection*{Your turn}
\begin{itemize}
\raggedright
  \item \emph{Say it:} \emph{thanganvũ}
  \item \emph{Say it:} \emph{thanganvũ}, once more
  \item \emph{Say it:} \emph{jo garbānũ sangīt śarū thaśe to bāḷko thanganśe.} — If the garba music starts, the children will be itching to dance.
  \item \emph{Recall:} say the Gujarati for to sing out, then the Gujarati for to chime, then say \emph{thanganvũ} again
\end{itemize}

\begin{tcolorbox}[breakable,colback=teal!4,colframe=teal!35!black,title={Before you move on}]
What does \gu{થનગનવું} mean? (\textquotedblleft{}To be itching to dance\textquotedblright{}.) Say it once more.
\end{tcolorbox}
